\section{The multiplicative group is cyclic}
\label{sec:finite-field-cyclic}

For a field $K$ with $q$ elements, its nonzero elements form a
commutative group $K^\times$ of order $N=q-1$. Lagrange gives
$a^N=1$ for every $a\ne0$. The stronger claim is that some element
has order exactly $N$.

Recall the polynomial root bound from interpolation: a nonzero polynomial
of degree $d$ over a field has at most $d$ roots. Indeed, division by
$X-a$ factors off a root $a$; another root $b\ne a$ must be a root of
the remaining factor, because $b-a$ can be cancelled. Induction proves
the bound. This cancellation is a field property.

\begin{lemma}[Combining coprime orders]
If commuting elements $a,b$ have relatively prime finite orders $r,s$,
then $ab$ has order $rs$.
\end{lemma}
\begin{proof}
Certainly $(ab)^{rs}=1$. If $(ab)^t=1$, put $z=a^t=b^{-t}$.
Then $z^r=z^s=1$. Choose integers $u,v$ with $ur+vs=1$; inverses
allow negative powers, so $z=(z^r)^u(z^s)^v=1$.
Consequently $r\mid t$ and $s\mid t$, hence $rs\mid t$.
Thus $rs$ is the least positive return exponent.
\end{proof}

\begin{theorem}[Cyclicity of finite-field multiplication]
\label{thm:finite-field-cyclic}
The multiplicative group of every finite field is cyclic.
\end{theorem}
\begin{proof}
If $N=1$, the group $\{1\}$ is cyclic. Otherwise, for each prime power
$\ell^r$ occurring exactly in $N$, the polynomial $X^{N/\ell}-1$
has fewer than $N$ roots. Choose $a\in K^\times$ with
$a^{N/\ell}\ne1$ and put $g_\ell=a^{N/\ell^r}$. Then
\[
g_\ell^{\ell^r}=a^N=1,\qquad
g_\ell^{\ell^{r-1}}=a^{N/\ell}\ne1.
\]
Its order divides $\ell^r$ but not $\ell^{r-1}$, so it is exactly
$\ell^r$. The product of these $g_\ell$ has order the product of their
pairwise coprime orders, by the lemma. This order is $N$, so its powers
exhaust $K^\times$.
\end{proof}

A generator of $K^\times$ is called a \emph{primitive element}.
It also generates $K$ as a field over its prime field, since every
nonzero element is a power of it. These two kinds of generation are
not equivalent in the opposite direction: in $\F_{49}$ the element
$s$ with $s^2=5$ generates the field over $\F_7$, but
$s^{12}=5^6=1$, so it does not generate the group of order $48$.

\begin{example}
In $\F_7^\times$, $6$ has order $2$ and $4$ has order $3$.
Their product $3$ therefore has order $6$; its powers are
$1,3,2,6,4,5$. A primitive element exists, but it is not a
distinguished choice of coordinate for the field.
\end{example}

\section{Existence and uniqueness of Galois fields}
\label{sec:finite-field-classification}

Two fields are \emph{isomorphic} if there is a bijection preserving
addition, multiplication, and $1$. Their elements may have different
names or polynomial presentations while obeying the same arithmetic.
We will prove that cardinality completely determines a finite field
up to isomorphism. A finite field is also called a \emph{Galois field},
and $\operatorname{GF}(q)$ is another notation for $\F_q$.

\subsection*{A polynomial construction supplies missing roots}

A nonconstant polynomial over $k$ is \emph{irreducible} if it cannot
be written as a product of two positive-degree polynomials over $k$.
Every nonconstant polynomial has an irreducible factor: if it is
reducible, factor it and continue with a factor of smaller positive
degree. The decreasing degrees make this process terminate.

\begin{lemma}[Adjoin one root]
\label{lem:adjoin-finite-root}
For a monic irreducible polynomial $h\in k[X]$ of degree $r$, the
polynomial residue ring $k[X]/(h)$ is a field containing $k$ and a
root of $h$. It has basis $1,t,\ldots,t^{r-1}$ over $k$, where $t$
is the class of $X$. If $|k|=q$, it has $q^r$ elements.
\end{lemma}
\begin{proof}
Division by the monic polynomial $h$ gives each residue a unique
representative of degree less than $r$. Constants embed, and $h(t)=0$.
If a representative $a$ is nonzero and has degree less than $r$,
irreducibility gives $\gcd(a,h)=1$. The polynomial Euclidean algorithm
works over a field by successive divisions, just as for integers.
Back substitution produces $u,v\in k[X]$ with $ua+vh=1$.
Reducing modulo $h$ gives an inverse for $a$. Thus the residue ring
is a field. Unique remainders give the stated basis and count.
\end{proof}

It follows that any polynomial over a finite field splits into linear
factors in some finite extension. To see the termination explicitly,
keep a factorization into known linear factors and a remaining factor.
If that factor is nonconstant, choose an irreducible factor of it and
adjoin a root by the lemma (or use a root already present). Divide out
the resulting linear factor and repeat. Each step lowers the remaining
degree by one, and every extension in the chain is finite. No prior
construction of an algebraic closure is required. A field generated by
all the roots of a polynomial is called its \emph{splitting field}.

\begin{theorem}[Classification of finite fields]
\label{thm:galois-field-uniqueness}
For every prime $p$ and integer $n\ge1$ there exists a field with
$p^n$ elements. Any two fields with that many elements are isomorphic
by an isomorphism fixing their identified prime fields $\F_p$.
Together with Theorem~\ref{thm:finite-field-size}, this describes all
possible finite-field cardinalities.
\end{theorem}
\begin{proof}
\emph{Existence.} Set $q=p^n$ and choose a finite extension $E$ of $\F_p$
in which $f(X)=X^q-X$ splits, using the preceding construction.
The formal derivative is defined by
$(\sum a_iX^i)'=\sum_{i\ge1}i a_iX^{i-1}$, with coefficients read in
the field; it obeys the product rule by direct multiplication.
Here $f'(X)=-1$, since $q=0$ in characteristic $p$.
A repeated root would make both $f$ and $f'$ vanish there: differentiate
$f=(X-a)^2h$ to see this. Thus $f$ has exactly $q$ distinct roots in $E$.

Let $S=\{a\in E:a^q=a\}$. The prime binomial identity, iterated $n$
times, gives $(a+b)^q=a^q+b^q$. Consequently $S$ is closed under
addition and additive inverses. It contains $0,1$, and for $a,b\in S$,
\[
(ab)^q=a^qb^q=ab,\qquad (a^{-1})^q=(a^q)^{-1}=a^{-1}\quad(a\ne0).
\]
So $S$ is a field, and it has exactly $q$ elements.

\emph{Uniqueness.} Let $K,L$ have $q$ elements. Their characteristics
are both $p$, by prime-power size and unique factorization. Identify
their prime fields. Lagrange in either multiplicative group gives
\[
X^q-X=\prod_{a\in K}(X-a)\quad\text{over }K,
\qquad
X^q-X=\prod_{b\in L}(X-b)\quad\text{over }L.
\]
Choose a generator $g$ of $K^\times$. There is a monic polynomial
$m\in\F_p[X]$ of least positive degree with $m(g)=0$, since
$g^q=g$. It is irreducible: a factorization into smaller positive
degrees would force one factor to vanish at $g$. Division with remainder
also shows that every polynomial vanishing at $g$ is divisible by $m$.
In particular $m\mid X^q-X$.

Since $X^q-X$ splits over $L$, $m$ has a root $b\in L$.
The evaluation map
\[
\F_p[X]/(m)\longrightarrow L,\qquad [a(X)]\longmapsto a(b)
\]
is well defined and injective: its domain is a field, and if a nonzero
element were sent to zero, multiplying by its inverse would send $1$
to zero as well. The same evaluation at
$g$ identifies $\F_p[X]/(m)$ with $\F_p[g]=K$, because all nonzero
elements of $K$ are powers of $g$. We have therefore embedded $K$
into $L$. Equal finite cardinalities make this embedding surjective,
giving the required isomorphism.
\end{proof}

\begin{warningbox}
Uniqueness means \emph{up to isomorphism}, not that there is only one
presentation or only one isomorphism. In
$K=\F_7(s)$, $s^2=5$, both $s\mapsto s$ and $s\mapsto-s$ preserve
the arithmetic. Likewise $\F_{p^n}$ is not the residue ring
$\Z/p^n\Z$ when $n>1$: their characteristics are $p$ and $p^n$.
\end{warningbox}

\begin{example}[Two presentations, one field]
Let $K=\F_3(s)$ with $s^2=2$. The polynomial $X^2-X-1$ is also
irreducible over $\F_3$, since its values at $0,1,2$ are $2,2,1$.
If $L=\F_3(t)$ with $t^2=t+1$, then
\[
(t+1)^2=t^2+2t+1=2.
\]
Thus $a+bs\mapsto a+b(t+1)$ is an isomorphism $K\to L$.
The alternative root $-(t+1)$ gives another isomorphism.
\end{example}

\section{Count matrices by choosing a basis}
\label{sec:finite-field-matrix-counts}

An invertible $d\times d$ matrix over $\F_q$ is the same thing as an
ordered basis of $\F_q^d$: its columns give the images of the standard
basis, and a basis gives unique solutions to every linear system.
After $r$ independent columns have been chosen, their span has $q^r$
vectors. The next column can be any of the $q^d-q^r$ vectors outside
that span. A full independent list of length $d$ is a basis by exchange.
Therefore, for $d\ge1$,
\[
\boxed{\left|\operatorname{GL}_d(\F_q)\right|
=\prod_{r=0}^{d-1}(q^d-q^r).}
\]
The factor for $r=0$ excludes the zero vector; order of columns matters.
For example, $|\operatorname{GL}_2(\F_2)|=(4-1)(4-2)=6$.

For $2\times2$ matrices, Appendix~\ref{app:matrices} defines
$\det\begin{psmallmatrix}a&b\\c&e\end{psmallmatrix}=ae-bc$ and
proves invertibility when this is nonzero. The matrices with determinant
$1$ form $\operatorname{SL}_2(\F_q)$. Every $a\in\F_q^\times$ occurs
as the determinant of $D_a=\operatorname{diag}(a,1)$, and
\[
S\longmapsto D_aS
\]
is a bijection from determinant $1$ matrices to determinant $a$
matrices, with inverse multiplication by $D_a^{-1}$.
There are $q-1$ disjoint fibers, all equally large. Thus
\[
\boxed{|\operatorname{SL}_2(\F_q)|
=\frac{(q^2-1)(q^2-q)}{q-1}=q(q^2-1).}
\]
The same fiber argument with the usual higher-dimensional determinant
gives $|\operatorname{SL}_d(\F_q)|=|\operatorname{GL}_d(\F_q)|/(q-1)$.
Neither matrix group is asserted to be cyclic; cyclicity above concerned
the nonzero \emph{field elements} under multiplication.
